\documentclass[11pt,a4paper]{article}
\usepackage[utf8]{inputenc}
\usepackage[T1]{fontenc}
\usepackage[french]{babel}
\usepackage{amsmath,amssymb,amsthm,geometry,hyperref}
\geometry{margin=2.4cm}
\newtheorem{proposition}{Proposition}
\title{Transport incrémental FE par espace auxiliaire\\Conservation, contraction et limite de petit pas}
\author{Prototype GPU SLFE, branche du fork keckj/rodin}
\date{16 novembre 2026}
\begin{document}\maketitle

\section{Construction}
Soient $V_h$ un espace DG $P_p$ sur triangles et $W_H$ un espace DG
cartésien $Q_q$. Les deux espaces contiennent la constante et représentent
le même domaine périodique. Leurs normes sont les normes $L^2$ exactes.
La projection orthogonale $R=P_{W_H}|_{V_h}$ a pour adjoint
$R^*=P_{V_h}|_{W_H}$. Les matrices sont
\[
 R=M_W^{-1}B,\qquad R^*=M_V^{-1}B^T,
 \qquad B_{ji}=\int_\Omega \psi_j\varphi_i.
\]
Le couplage $B$ est statique : il ne dépend ni de la vitesse, ni du pas.
Le calcul des intersections triangle--rectangle est effectué une fois.
Les intersections mobiles de cellules transportées ne sont pas nécessaires
à cette construction. La précision des intégrales statiques reste essentielle.

Soit $S_{\Delta t}:W_H\to W_H$ un transport conservatif, préservant la
constante et contractant dans la norme $M_W$. Le transfert proposé est
\[
 \boxed{T_{\Delta t}=I+R^*(S_{\Delta t}-I)R.}
\]
Il projette le champ vers le support auxiliaire, y calcule son évolution,
puis retourne seulement l'incrément. Une projection du champ entier
donnerait $R^*S_{\Delta t}R$ et introduirait $R^*R$ même sans déplacement.

\section{Propriétés discrètes}
\begin{proposition}
Si $R$ et $S$ satisfont les hypothèses ci-dessus, $T$ conserve la masse,
préserve la constante et vérifie $\|Tv\|_{M_V}\leq\|v\|_{M_V}$.
Si $S_0=I$, alors $T_0=I$.
\end{proposition}
\begin{proof}
Pour la constante $1_V$, $R1_V=1_W$, d'où $T1_V=1_V$.
En notant $d=(S-I)Rv$, la conservation dans $W_H$ et l'adjonction donnent
\[
 \langle1_V,Tv-v\rangle_V=\langle R1_V,d\rangle_W
 =\langle1_W,(S-I)Rv\rangle_W=0.
\]
Comme $\|R^*d\|_V\leq\|d\|_W$,
\begin{align*}
 \|Tv\|_V^2
 &=\|v\|_V^2+2\langle Rv,d\rangle_W+\|R^*d\|_V^2\\
 &\leq\|v\|_V^2+\|SRv\|_W^2-\|Rv\|_W^2
 \leq\|v\|_V^2.
\end{align*}
L'identité à pas nul résulte directement de la formule de $T$.
\end{proof}

La dissipation se décompose exactement en
\[
 \|v\|_V^2-\|Tv\|_V^2
 =\underbrace{\|Rv\|_W^2-\|SRv\|_W^2}_{\text{transport auxiliaire}}
 +\underbrace{\|d\|_W^2-\|R^*d\|_V^2}_{\text{retour de l'incrément}}.
\]
Cette preuve s'applique au transfert linéaire entier et donc à la différence
de deux champs. Aucun limiteur dépendant du champ ni ajustement de moments
positifs n'est requis. Elle n'implique pas la positivité des valeurs du champ.

Si $S_{\Delta t}=I+\Delta t L_W+o(\Delta t)$, alors
\[
 T_{\Delta t}=I+\Delta t R^*L_WR+o(\Delta t).
\]
Il n'existe donc pas de dissipation fixe due à une double projection à
chaque pas lorsque $\Delta t\to0$. Le générateur obtenu reste celui d'un
problème spatial approché ; l'identité à pas nul ne prouve pas sa consistance.

\section{Réalisation sur GPU}
La base triangulaire est modale et orthonormale en moyenne. Le support
cartésien utilise les noeuds de Gauss tensorisés, dont la matrice de masse
est exactement diagonale pour $Q_q$. Le produit croisé est un polynôme de
degré total au plus $p+2q$ sur chaque intersection affine. Une quadrature
de ce degré fournit les matrices statiques à l'arrondi près.

Pour la translation, $S$ est composé de deux projections DG 1D exactes.
Chaque direction ne couple que deux cellules voisines du pied décalé par
la partie entière du déplacement. Le nombre de blocs par ligne ne dépend
donc pas du nombre de cellules franchies.

Pour le cisaillement $u=(a\sin(2\pi y),0)$, chaque ligne de noeuds de Gauss
en $y$ possède sa propre translation en $x$. Chaque ligne est contractante
dans sa norme de Gauss ; leur somme pondérée l'est également. L'énergie
discrète correspond à la norme $L^2$ exacte du polynôme $Q_q$ reconstruit.
La variation non polynomiale en $y$ introduit une erreur de consistance
distincte de cette propriété de stabilité.

Le code ne fournit pas encore de transport auxiliaire général pour une
vitesse FE non séparable. Une décomposition arbitraire en directions ne
doit pas être supposée contractante pour un flot incompressible quelconque.

\section{Résolution auxiliaire et limites}
Si $v\in\ker R$, alors $Tv=v$ : les modes invisibles au support auxiliaire
restent immobiles. Pour $q=p$ et un carré auxiliaire par carré logique de
deux triangles, les dimensions sont $\dim V_h=n^2(p+1)(p+2)$ et
$\dim W_H=n^2(p+1)^2$. Un noyau est donc inévitable. Les expériences
rapportent cette dimension et comparent des supports plus fins. Le rang
complet observé sur un support enrichi ne constitue pas une garantie pour
toutes les géométries.

L'ordre spatial, l'erreur de phase et l'accumulation temporelle sont mesurés
séparément. Une preuve uniforme de convergence, une stratégie de choix de
$H,q$, les frontières non périodiques, le transport compressible, un support
auxiliaire adapté à des géométries complexes et le coût MPI restent ouverts.

Le bilan global de masse provient du transfert de l'incrément conservatif.
Il ne suffit pas à identifier un flux physique unique sur chaque face du
maillage FE d'origine ; cette propriété locale nécessite une analyse distincte.

\section{Positionnement}
Le transfert d'incréments possède des antécédents explicites dans les
méthodes FLIP. Les projections adjointes et les espaces auxiliaires sont
également des outils établis. La contribution éventuelle ne peut donc pas
être l'idée générale de retourner une variation. La formulation FE avec
projections $L^2$ statiques, le transport DG auxiliaire contractant, le
traitement des modes invisibles, les bornes de consistance et les coûts GPU
doivent être comparés précisément aux travaux antérieurs.

\begin{thebibliography}{9}
\bibitem{flip} J. U. Brackbill et H. M. Ruppel,
\emph{FLIP: A method for adaptively zoned, particle-in-cell calculations
of fluid flows in two dimensions}, J. Comput. Phys. 65 (1986), 314--343.
\bibitem{assign} D. Duque et P. Español,
\emph{Assignment of fields from particles to mesh},
\url{https://arxiv.org/abs/1610.05258}.
\bibitem{coflip} \emph{Fluid Implicit Particles on Coadjoint Orbits},
\url{https://arxiv.org/abs/2406.01936}.
\end{thebibliography}
\end{document}
